\documentclass[10pt,a4]{examenEPN0}
\usepackage[T1]{fontenc}
\usepackage{comandosEPN}

\usepackage{empheq}

\newcommand{\itemcp}{\stepcounter{enumi}\item[ \textcolor{red}{\theenumi.$^*$}] }

\newcommand{\iitemcp}{\stepcounter{enumii}\item[ \textcolor{red}{\theenumii)$^*$}] }


\materia{Álgebra Lineal}
%\autor{Departamento de Formación Básica}
\autor{GR \hspace{1cm}.}
%\fecha{Semestre 2026-A}
\fecha{Nombre: \hspace{9cm} Cédula:}
\logouno[1.6cm]{Logo_EPN.eps}
\logodos[2.6cm]{Logo_DFB.eps}

\newcommand{\la}{\lambda}

%Datos
\newcommand{\datosAlgebra}[2]{
\EnExamen
{
\vspace{5mm}
    \noindent
    Nombres: \hrulefill\
    Carrera: \rule{4.5cm}{0.4pt}\
    \ifthenelse{\equal{#1}{}}%
    {Paralelo: \underline{\parbox{1.8cm}{\parbox{1cm}{}}}\\[3mm]}
    {Paralelo: \underline{\parbox{1.8cm}{\centering #1}}\\[3mm]}
    C\'edula: \rule{5cm}{0.4pt}\
    Firma: \hrulefill\
    Fecha: \underline{\parbox{2.1cm}{\centering #2}}
}}


% \newcommand{\indicacionesPrueba}{%
% \EnExamen{
% \begin{indicaciones}
%     \item La prueba tiene una duraci\'on  de \textbf{120 minutos.}
%     \item No se permite la utilizaci\'on de dispositivos electr\'onicos (celular, tablets, calculadora, etc.)
%     \item Leer detenidamente los enunciados, \textbf{justificar la respuesta para acceder al puntaje total}.
%     \item La prueba es completamente individual. Los estudiantes deberán dejar sus teléfonos celulares guardados y no podrán tenerlos consigo (por ejemplo, en los bolsillos de su ropa) durante toda la evaluación. En caso de detectarse algún acto de deshonestidad académica o el incumplimiento de esta norma, no podrán seguir rindiendo la prueba y su calificación será \textbf{cero}.
% \end{indicaciones}
% \pagestyle{empty}
% }}

\newcommand{\indicacionesPrueba}{%
\EnExamen{
\begin{indicaciones}
    \item La prueba tiene una duraci\'on  de \textbf{60 minutos.}
    \item Leer detenidamente los enunciados, \textbf{justificar la respuesta} para acceder al puntaje total.
    \item No se permite la utilizaci\'on de dispositivos electr\'onicos (celular, tablets, calculadora, etc.)
    \item La prueba es completamente individual
    \item No se permite portar teléfonos celulares, relojes inteligentes ni otros dispositivos electrónicos durante el examen, aun cuando permanezcan en los bolsillos. Estos dispositivos deberán guardarse en la mochila antes del inicio de la evaluación y esta deberá mantenerse cerrada durante todo el examen.
    \item He leído las indicaciones y acepto que el incumplimiento de estas se penaliza con \textbf{cero}. Firma:...................................

    
\end{indicaciones}
\pagestyle{empty}
}}

\newcommand{\indicacionesExamen}{
\EnExamen{\begin{indicaciones}
    \item La evaluación tiene una duraci\'on  de \textbf{120 minutos.}
    \item No se permite la utilizaci\'on de dispositivos electr\'onicos (celular, tablets, calculadora, etc.)
    %\item Leer detenidamente los enunciados, justificar la respuesta para acceder al puntaje total.
    \item El examen es completamente individual. En caso de detectarse alg\'un acto de deshonestidad acad\'emica no podr\'a seguir rindiendo la prueba y su calificaci\'on ser\'a \textbf{cero}.
\end{indicaciones}
}\pagestyle{empty}
}

%% Imprimir las una hoja A4, dos veces en A5
\newcommand{\imprimirDoble}
{%
\usepackage{pgfmorepages}
\pgfmorepagesloadextralayouts
\pgfpagesuselayout{repeated 2-up}[a4paper,landscape]
\thispagestyle{empty}
}


%Este sí se puede comentar
\newenvironment{solucion}{\color{blue}\proof[Solución]}{\endproof\color{black}}




% \renewcommand{\@autor}{
% ---\\
% \phantom{.} 
% ...}
% \makeatother

\examen{Prueba 01: solucionario}
\makeatletter
\renewcommand{\@autor}{Departamento de Formación Básica}
\renewcommand{\@fecha}{GR5M -- Miércoles de 7 a 9 -- 21/10/2026}
\makeatother

\begin{document}
\encabezado

\section*{1. Sistema homogéneo (5 puntos)}
La matriz de coeficientes tiene dos columnas, de modo que hay dos
incógnitas, $x_1$ y $x_2$. Reducimos por filas:
\[
    A=\begin{pmatrix}1&0\\1&1\\0&1\end{pmatrix}
    \xrightarrow{F_2\leftarrow F_2-F_1}
    \begin{pmatrix}1&0\\0&1\\0&1\end{pmatrix}
    \xrightarrow{F_3\leftarrow F_3-F_2}
    \begin{pmatrix}1&0\\0&1\\0&0\end{pmatrix}.
\]
Hay dos pivotes, por lo que $\rang(A)=2$. Como el sistema es homogéneo,
es compatible y $\rang(A\mid\mathbf{0})=\rang(A)=2$. El rango coincide
con el número de incógnitas; por tanto, no hay variables libres y la
solución es única. Las ecuaciones reducidas dan $x_1=0$ y $x_2=0$.
El conjunto solución es
\[
    \boxed{S=\{(0,0)\}.}
\]

\section*{2. Potencia de un producto de matrices (5 puntos)}
La demostración propuesta es incorrecta. El paso 1 usa la definición
de cuadrado y el paso 2 usa la asociatividad del producto matricial.
El error está en el \textbf{paso 3}: se reemplaza $BA$ por $AB$ sin
justificación. En general, el producto de matrices no es conmutativo.
El paso 4 sería válido por asociatividad si se hubiera llegado
correctamente a $A(AB)B$.

Para refutar la afirmación, tomemos
\[
    A=\begin{pmatrix}1&0\\0&0\end{pmatrix},\qquad
    B=\begin{pmatrix}0&1\\1&0\end{pmatrix}.
\]
Entonces
\[
    AB=\begin{pmatrix}0&1\\0&0\end{pmatrix},\qquad
    (AB)^2=\begin{pmatrix}0&0\\0&0\end{pmatrix}.
\]
Por otra parte,
\[
    A^2=\begin{pmatrix}1&0\\0&0\end{pmatrix}=A,\qquad
    B^2=\begin{pmatrix}1&0\\0&1\end{pmatrix}=I_2,
\]
y, en consecuencia,
\[
    A^2B^2=A=\begin{pmatrix}1&0\\0&0\end{pmatrix}
    \ne\begin{pmatrix}0&0\\0&0\end{pmatrix}=(AB)^2.
\]
Un solo contraejemplo basta para refutar una afirmación universal.
Si se añade la hipótesis $AB=BA$, el razonamiento sí demuestra la igualdad.

\newpage
\section*{3. Sistema con parámetro (5 puntos)}
\noindent\textbf{a) Matriz ampliada y operaciones indicadas (1 punto).}
La matriz ampliada es
\[
    (A\mid b)=\left(\begin{array}{ccc|c}
        1&1&1&1\\0&1&k&2\\1&2&2&3
    \end{array}\right).
\]
Aplicamos las operaciones en el orden solicitado:
\[
    \left(\begin{array}{ccc|c}
        1&1&1&1\\0&1&k&2\\1&2&2&3
    \end{array}\right)
    \xrightarrow{F_3\leftarrow F_3-F_1}
    \left(\begin{array}{ccc|c}
        1&1&1&1\\0&1&k&2\\0&1&1&2
    \end{array}\right)
    \xrightarrow{F_3\leftarrow F_3-F_2}
    \left(\begin{array}{ccc|c}
        1&1&1&1\\0&1&k&2\\0&0&1-k&0
    \end{array}\right).
\]
En la segunda operación se utilizó la tercera fila ya modificada.

\noindent\textbf{b) Rangos y clasificación (4 puntos).}
Por el teorema de Rouché--Frobenius, comparamos los rangos entre sí
y con el número de incógnitas, que es 3:
\begin{center}
    \begin{tabular}{c|c|c|l}
        Condición & $\rang(A)$ & $\rang(A\mid b)$ & Tipo de solución\\
        \hline
        $k\ne1$ & 3 & 3 & Única\\
        $k=1$ & 2 & 2 & Infinitas soluciones
    \end{tabular}
\end{center}
\begin{itemize}
    \item Si $k\ne1$, la última fila tiene pivote. Hay tres pivotes en
          ambas matrices y el rango coincide con el número de incógnitas.
          La última ecuación da $x_3=0$; luego $x_2=2$ y $x_1=-1$.
          Por tanto, $S=\{(-1,2,0)\}$.
    \item Si $k=1$, la última fila es nula en la matriz ampliada.
          Los rangos son iguales a 2, menores que 3, por lo que queda una
          variable libre. Al tomar $x_3=t$, obtenemos $x_2=2-t$ y
          $x_1=1-(2-t)-t=-1$. Así,
          \[
              S=\{(-1,2-t,t):t\in\R\}.
          \]
\end{itemize}
\textbf{No existe ningún valor de $k$ que produzca un sistema incompatible:}
la última fila nunca tiene la forma $[\,0\;0\;0\mid c\,]$ con $c\ne0$.

\section*{4. Modelo de mezclas (5 puntos)}
\noindent\textbf{a) Incógnitas y unidades (1 punto).}
Para $i=1,2,3$, sea $x_i$ el volumen de la mezcla preparada $M_i$
que se utiliza, medido en metros cúbicos. Por su interpretación física,
$x_1,x_2,x_3\ge0$.

\noindent\textbf{b) Ecuaciones de los materiales (3 puntos).}
Cada proporción multiplica el volumen de la mezcla correspondiente:
\[
    \left\{\begin{aligned}
        0.30x_1+0.20x_2+0.40x_3&=28 &&\text{(cemento)},\\
        0.40x_1+0.50x_2+0.30x_3&=42 &&\text{(arena)},\\
        0.30x_1+0.30x_2+0.30x_3&=30 &&\text{(grava)}.
    \end{aligned}\right.
\]
Ambos miembros de cada ecuación se expresan en metros cúbicos.
Al sumar las tres ecuaciones se obtiene $x_1+x_2+x_3=100$, por lo que
el volumen total solicitado ya está incluido en el modelo.

\noindent\textbf{c) Forma matricial (1 punto).}
El sistema es $Ax=b$, donde
\[
    \boxed{
    A=\begin{pmatrix}
        0.30&0.20&0.40\\
        0.40&0.50&0.30\\
        0.30&0.30&0.30
    \end{pmatrix},\qquad
    x=\begin{pmatrix}x_1\\x_2\\x_3\end{pmatrix},\qquad
    b=\begin{pmatrix}28\\42\\30\end{pmatrix}.}
\]
Las filas de $A$ corresponden a cemento, arena y grava; sus columnas,
a $M_1$, $M_2$ y $M_3$, respectivamente. El planteamiento se completa con
la restricción $x_1,x_2,x_3\ge0$. No se requiere resolver el sistema.

\end{document}
